\begin{figure}[tbp]
\centering
\begin{minipage}[t]{.40\linewidth}
\centering\ViewHeading{(a) One common denominator}
\begin{tikzpicture}[scale=.8]
\draw[viewgray!50,step=1] (0,0) grid (4,4);
\foreach \horizontal in {0,...,4}
 \foreach \vertical in {0,...,4} {
 \draw[viewblue!50,fill=viewblue!8] (\horizontal-.15,\vertical-.15) rectangle (\horizontal+.15,\vertical+.15);
 \fill[viewblue] (\horizontal,\vertical) circle (1.6pt);
 }
\node[vseed,label=above right:$x$] at (2.12,1.1) {};
\node[font=\footnotesize,align=center] at (2,-.8)
 {$p/q$, schematic $q=4$ grid\\neighborhood radius $cq^{-(1+\alpha)}$};
\end{tikzpicture}
\end{minipage}\hfill
\begin{minipage}[t]{.57\linewidth}
\centering\ViewHeading{(b) Comparison, not equality of loci}
\begin{tikzpicture}
\node[vbox,text width=6.8cm] (bound) at (0,0)
 {$4\Delta_q(x)\leq b_q(x)\leq2\pi\sqrt d\,\Delta_q(x)$};
\node[vresult,text width=6.8cm] (sandwich) at (0,-1.7)
 {$\displaystyle W_{\alpha,C/(2\pi\sqrt d)}\subseteq R_{\alpha,C}\subseteq W_{\alpha,C/4}$};
\draw[varrow] (bound)--(sandwich);
\node[vinput,text width=6.8cm] at (0,-3.55)
 {Classical dimension theorem\\and constant comparison\\
 $\displaystyle\dim_{\rm H}R_{\alpha,C}=\frac{d+1}{1+\alpha}$\\only in the stated range $\alpha>1/d$};
\end{tikzpicture}
\end{minipage}
\caption{The same integer denominator is used in every coordinate.
Panel (a) is a finite grid schematic, not a picture of the limsup set.
The two comparison constants in (b) are essential: equality of the
resulting Hausdorff dimensions does not identify the selected sets at
the same tolerance.}
\label{fig:arithmetic-recurrence}
\end{figure}